\exposetitle{II}{Tangent bundles. Lie algebras}
\label{II}

\begin{center}
\textit{by M.~Demazure}
\end{center}

{\renewcommand{\thefootnote}{(0)}\nde{Version of 14/10/2024}}%
We propose in this expos\'e to construct the analogue, in the theory of
schemes, of the \emph{tangent bundles} and \emph{Lie algebras} of the
classical theory. It will however be useful not to restrict oneself to
schemes properly so called, but also to take an interest in certain functors
on the category of schemes which are not necessarily \emph{representable}
(for example functors \(\Hom\), \(\Norm\), etc.). As was announced in the
previous expos\'e (\cf I~1.1), we shall identify a scheme with the functor
associated with it.

On the other hand, the constructions set out below go beyond the framework
of the theory of schemes. They are equally valid, for example, in the theory
of \emph{analytic spaces} with nilpotent elements, modulo a few modifications
of detail.

Before beginning this construction, we must set down a few general
definitions which complete those of~I~1.7.

\sgasection{1}{The functors \(\Hom_{Z/S}(X,Y)\)}
\label{II.sec.1}

Let us take up again the notations of~I~1.1. One identifies the category
\(\mathcal{C}\) with a full subcategory of
\(\widehat{\mathcal{C}}=\Hom(\mathcal{C}^{\circ},\Ens)\) (in particular one
omits the underlinings%
{\renewcommand{\thefootnote}{(1)}\nde{Rendered in this edition by bold
characters \(\mathbf{F}\), \(\mathbf{G}\), \(\mathbf{V}\), \(\mathbf{W}\),
\cf Expos\'e~I. However, for functors such as \(\uNorm\) and \(\uCentr\)
(\cf~5.2) the underlinings have been kept as in the original, and one has
added them for the functor \(\Lie\), in order to distinguish the functor
\(\uLie(G/S)\) from the \(\Gamma(S,\OO_S)\)-Lie algebra
\(\Lie(G/S)=\uLie(G/S)(S)\), used for example in
Expos\'e~$\mathrm{VII}_{\mathrm{A}}$.}}
which allowed us to distinguish graphically an object of
\(\widehat{\mathcal{C}}\) from an object of \(\mathcal{C}\)).

Let us consider the following situation: four objects of
\(\widehat{\mathcal{C}}\), denoted \(S\), \(X\), \(Y\), \(Z\), the first
being in fact an object of \(\mathcal{C}\), \(X\) and \(Y\) over \(Z\), \(Z\)
over \(S\):
\[
\xymatrix@C=2.2em@R=1.4em{
X \ar[dr]_{p_X} & & Y \ar[dl]^{p_Y} \\
& Z \ar[d] & \\
& S\,. &
}
\]

\begin{definition}{1.1}
\label{II.1.1}
One defines an object of \(\widehat{\mathcal{C}}_{/S}\), denoted
\(\Hom_{Z/S}(X,Y)\), by:
\begin{equation}
\Hom_{Z/S}(X,Y)(S')
=\Hom_{Z_{S'}}(X_{S'},Y_{S'})
=\Hom_Z(X\times_S S',Y),
\tag{1}
\end{equation}
for every object \(S'\) of \(\mathcal{C}_{/S}\). One sees at once that
\(\Hom_{Z/S}(X,Y)\) is nothing other than the subobject of
\(\Hom_S(X,Y)\) formed of the morphisms compatible with \(p_X\) and \(p_Y\),
that is, the kernel of the pair of morphisms
\[
\xymatrix{
\Hom_S(X,Y) \ar@<0.5ex>[r] \ar@<-0.5ex>[r] & \Hom_S(X,Z)
}
\]
defined, the first by composition with \(p_Y\), the second as being the
constant morphism whose ``image'' is \(p_X\).

{\renewcommand{\thefootnote}{(2)}\nde{Points (2), (3), (4) have been spelled
out; in particular, (4) will be used in~3.11.}}%
On the other hand, one sees as in~I~1.7 that, for every object \(T\) of
\(\widehat{\mathcal{C}}\) over \(S\), one has a natural bijection:
\begin{equation}
\Hom_S\bigl(T,\Hom_{Z/S}(X,Y)\bigr)
\simeq
\Hom_Z(X\times_S T,Y).
\tag{2}
\end{equation}
Moreover, by~I~1.7.1, if \(E\), \(F\) are objects of \(\widehat{\mathcal{C}}\)
over \(Z\), one has:
\[
\Hom_Z\bigl(E,\Hom_Z(F,Y)\bigr)
\simeq
\Hom_Z(E\times_Z F,Y)
\simeq
\Hom_Z\bigl(F,\Hom_Z(E,Y)\bigr).
\]
Applying this to \(E=X\) and \(F=Z\times_S T\), one obtains natural
bijections, for every object \(T\) of \(\widehat{\mathcal{C}}_{/S}\):
\begin{equation}
\Hom_S\bigl(T,\Hom_{Z/S}(X,Y)\bigr)
\simeq
\Hom_Z(X\times_S T,Y)
\simeq
\begin{cases}
\Hom_Z\bigl(Z\times_S T,\Hom_Z(X,Y)\bigr)\\
\Hom_Z\bigl(X,\Hom_Z(Z\times_S T,Y)\bigr).
\end{cases}
\tag{3}
\end{equation}
Moreover, these bijections are functorial in \(T\), so one obtains
isomorphisms of \(S\)-functors:
\begin{equation}
\vcenter{\xymatrix@C=1.4em@R=1.6em{
\Hom_S\bigl(T,\Hom_{Z/S}(X,Y)\bigr)
  \ar[rr]^{\sim} \ar[dr]_{\simeq}
&&
\Hom_{Z/S}\bigl(X,\Hom_Z(Z\times_S T,Y)\bigr)
  \ar[dl]^{\simeq}
\\
& \Hom_{Z/S}(X\times_S T,Y) &
}}
\tag{4}
\end{equation}

Let us point out two particular cases of the definition. If \(Z=S\), one
has:
\[
\Hom_{S/S}(X,Y)=\Hom_S(X,Y).
\]
On the other hand, when \(X=Z\), one sets
\begin{equation}
\prod_{Z/S}Y=\Hom_{Z/S}(Z,Y),
\tag{5}
\end{equation}
one therefore has by definition
\[
\Bigl(\prod_{Z/S}Y\Bigr)(S')
=\Hom_Z(Z\times_S S',Y)
\simeq\Gamma(Y_{S'}/Z_{S'}).
\]
The functor
\(\prod_{Z/S}\colon\widehat{\mathcal{C}}_{/Z}\to\widehat{\mathcal{C}}_{/S}\)
is right adjoint to the functor of base change from \(S\) to \(Z\): for every
\(S\)-functor \(U\) one has
\[
\Hom_S\Bigl(U,\prod_{Z/S}Y\Bigr)
=\Hom_Z(U\times_S Z,Y).
\]
(If \(\mathcal{C}=\Sch\) and if \(Z\) is an \(S\)-scheme, the functor
\(\prod_{Z/S}\) is called ``\emph{Weil restriction of scalars}''.)%
{\renewcommand{\thefootnote}{(3)}\nde{The two preceding sentences have been
added.}}
Let us also note that one has an isomorphism:
\begin{equation}
\Hom_{Z/S}(X,Y)
\simeq
\Hom_{X/S}(X,Y\times_Z X)
=\prod_{X/S}(Y\times_Z X),
\tag{6}
\end{equation}
which gives in particular, for \(Z=S\), an isomorphism:
\begin{equation}
\Hom_S(X,Y)\simeq\prod_{X/S}Y_X.
\tag{7}
\end{equation}
\end{definition}

\begin{remark}{1.2}
\label{II.1.2}
The functor \(Y\mapsto\Hom_{Z/S}(X,Y)\) commutes with the product in the
following sense: one has a functorial isomorphism
\begin{equation}
\Hom_{Z/S}(X,Y\times_Z Y')
\simeq
\Hom_{Z/S}(X,Y)\times_S\Hom_{Z/S}(X,Y').
\tag{*}
\end{equation}
It follows that if \(Y\) is a \(Z\)-group, \resp a \(Z\)-ring, etc., then
\(\Hom_{Z/S}(X,Y)\) is an \(S\)-group, \resp an \(S\)-ring, etc.
\end{remark}

\begin{remark}{1.3}
\label{II.1.3}
{\renewcommand{\thefootnote}{(4)}\nde{This remark has been added; it is
useful for Corollary~3.11.1.}}%
Moreover, let \(\pi\colon M\to Y\) be a \(Y\)-functor in
\(\mathbf{O}_Y\)-modules (\cf I,~4.3.3.1). Set
\(H=\Hom_{Z/S}(X,Y)\). Then \(\Hom_{Z/S}(X,M)\) is endowed with a natural
structure of \(\mathbf{O}_H\)-module; more precisely, for every \(H'\) over
\(H\), \(\Hom_H(H',\Hom_{Z/S}(X,M))\) is endowed with a natural structure of
\(\mathbf{O}(H'\times_S X)\)-module.

Indeed, let us denote by \(m\colon M\times_Y M\to M\) and
\(\lambda\colon\mathbf{O}_Y\times_Y M\to M\) the morphisms defining the
structures of \(Y\)-group (abelian) and of \(\mathbf{O}_Y\)-module. Let
\(H'\) be an \(S\)-scheme over \(H=\Hom_{Z/S}(X,Y)\), that is, one is given a
\(Z\)-morphism \(f\colon X\times_S H'\to Y\), which therefore makes
\(X\times_S H'\) a \(Y\)-object. Then
\[
\Hom_H\bigl(H',\Hom_{Z/S}(X,M)\bigr)
\]
is the set of \(Z\)-morphisms \(\phi\colon X\times_S H'\to M\) such that
\(\pi\circ\phi=f\), that is, of \(Y\)-morphisms \(X\times_S H'\to M\).

If \(\phi\), \(\psi\) are two such morphisms, one defines \(\phi+\psi\) as the
composite \(Y\)-morphism
\[
\xymatrix{
X\times_S H' \ar[r]^{\phi\times\psi} & M\times_Y M \ar[r]^{m} & M
}
\]
and one checks that this endows \(\Hom_{Z/S}(X,M)\) with the structure of an
abelian group over \(H=\Hom_{Z/S}(X,Y)\).%
{\renewcommand{\thefootnote}{(5)}\nde{Of course, if \(M\) is a \(Y\)-group
(not necessarily abelian) and if one sets \(\phi\cdot\psi=m\circ(\phi\times\psi)\),
this makes \(\Hom_{Z/S}(X,M)\) a group over \(\Hom_{Z/S}(X,Y)\).}}

Likewise, if \(a\) is an element of \(\mathbf{O}(X\times_S H')\), \ie an
\(S\)-morphism \(a\colon X\times_S H'\to\mathbf{O}_S\), one defines \(a\phi\)
as the composite \(\lambda\circ(a\times\phi)\), where \(a\times\phi\) denotes
the \(Y\)-morphism from \(X\times_S H'\) to
\(\mathbf{O}_Y\times_Y M\simeq\mathbf{O}_S\times_S M\) with components \(a\)
and \(\phi\); one checks that this endows
\(\Hom_H(H',\Hom_{Z/S}(X,M))\) with the structure of an
\(\mathbf{O}(X\times_S H')\)-module, functorial in the \(H\)-object \(H'\).
\end{remark}

\sgasection{2}{The schemes \(\mathrm{I}_S(\mathcal{M})\)}
\label{II.sec.2}

\begin{definition}{2.1}
\label{II.2.1}
Let \(S\) be a scheme and \(\mathcal{M}\) a quasi-coherent
\(\OO_S\)-module. One denotes by \(\mathcal{D}_{\OO_S}(\mathcal{M})\) the
quasi-coherent \(\OO_S\)-algebra \(\OO_S\oplus\mathcal{M}\)
(\emph{where \(\mathcal{M}\) is regarded as a square-zero ideal}). One
denotes by \(\mathrm{I}_S(\mathcal{M})\) the \(S\)-scheme
\(\Spec\mathcal{D}_{\OO_S}(\mathcal{M})\).%
{\renewcommand{\thefootnote}{(6)}\nde{Note that \(\mathrm{I}_S(\mathcal{M})\)
has the same underlying space as \(S\).}}
In particular one denotes
\(\mathcal{D}_{\OO_S}=\mathcal{D}_{\OO_S}(\OO_S)\),
\(\mathrm{I}_S=\mathrm{I}_S(\OO_S)\) and one calls them respectively the
\emph{algebra of dual numbers over \(S\)} and the \emph{scheme of dual
numbers over \(S\)}.

Then \(\mathcal{M}\mapsto\mathrm{I}_S(\mathcal{M})\) is a contravariant
functor from the category of quasi-coherent \(\OO_S\)-modules to that of
\(S\)-schemes. In particular the morphisms \(0\to\mathcal{M}\) and
\(\mathcal{M}\to 0\) define respectively the structural morphism
\(\rho\colon\mathrm{I}_S(\mathcal{M})\to\mathrm{I}_S(0)=S\) and a section
\(\varepsilon_{\mathcal{M}}\) of the latter, which one calls the
\emph{zero section}.%
{\renewcommand{\thefootnote}{(7)}\nde{Compare with~3.1.1.}}
\end{definition}

\numpara{2.1.1}\label{II.2.1.1}
{\renewcommand{\thefootnote}{(8)}\nde{The original has been set out in detail
in what follows, and the numbering 2.1.1 to 2.1.3 has been added.}}%
As \(\mathcal{M}\to\mathrm{I}_S(\mathcal{M})\) is a contravariant functor,
% typo?: printed \to, not \mapsto; kept as printed
every \(a\in\End_{\OO_S}(\mathcal{M})\) defines an \(S\)-endomorphism \(a^*\)
of \(\mathrm{I}_S(\mathcal{M})\), and one has \(1^*=\mathrm{id}\),
\((ab)^*=b^*\circ a^*\), \(0^*=\varepsilon_{\mathcal{M}}\circ\rho\) and
\(a^*\circ\varepsilon_{\mathcal{M}}=\varepsilon_{\mathcal{M}}\).
Consequently, the \(S\)-scheme \(\mathrm{I}_S(\mathcal{M})\) is endowed with a
right action of the multiplicative monoid of \(\End_{\OO_S}(\mathcal{M})\),
which commutes with the \(S\)-morphisms
\(\mathrm{I}_S(\mathcal{M})\to\mathrm{I}_S(\mathcal{M}')\) coming from
morphisms \(\mathcal{M}'\to\mathcal{M}\); in particular, the operators
\(a^*\) preserve the zero section of \(\mathrm{I}_S(\mathcal{M})\).

For every \(a\in\End_{\OO_S}(\mathcal{M})\) and \(f\colon S'\to S\) and
\(m\in\mathrm{I}_S(\mathcal{M})(S')\), one will denote
\(m\cdot a=a^*(m)\); then \(m\cdot 1=m\),
\((m\cdot a)\cdot b=m\cdot(ab)\),
\(m\cdot 0=\varepsilon_{\mathcal{M}}(\rho(m))\) and, if
\(m=\varepsilon_{\mathcal{M}}(f)\), then \(m\cdot a=m\).%
{\renewcommand{\thefootnote}{(9)}\nde{In the sequel, one will be concerned
mainly with the case where \(a\in\mathbf{O}(S)\), acting by homotheties by
\(\mathcal{M}\). For example, if \(\mathcal{M}\) is a free \(\OO_S\)-module
of rank \(r\), then, for every \(S'\to S\),
\(\Hom_S(S',\mathrm{I}_S(\mathcal{M}))\) is identified with the set of
\(r\)-tuples \((\varepsilon_1,\ldots,\varepsilon_r)\in\mathbf{O}(S')\) such
that \(\varepsilon_i\varepsilon_j=0\) for all \(i\), \(j\), and one has
\((\varepsilon_1,\ldots,\varepsilon_r)\cdot a
=(a\varepsilon_1,\ldots,a\varepsilon_r)\) for every
\(a\in\mathbf{O}(S)\).}}

\begin{remark}{2.1.2}
\label{II.2.1.2}
The formation of the \(\mathrm{I}_S(\mathcal{M})\) \emph{commutes with
extension of the base}: one has canonical isomorphisms
\[
\mathrm{I}_S(\mathcal{M})_{S'}
\simeq
\mathrm{I}_{S'}(\mathcal{M}\otimes_{\OO_S}\OO_{S'}).
\]
For simplicity, one will denote
\(\mathrm{I}_{S'}(\mathcal{M})=\mathrm{I}_S(\mathcal{M})_{S'}\); more
generally, if \(X\) is an \(S\)-functor (not necessarily representable), one
will denote \(\mathrm{I}_X(\mathcal{M})=\mathrm{I}_S(\mathcal{M})\times_S X\).
\end{remark}

\numpara{2.1.3}\label{II.2.1.3}
{\renewcommand{\thefootnote}{(10)}\nde{This paragraph has been added; it will
be useful in~3.4.2 and in~4.6.2.}}%
By what precedes, the multiplicative monoid of \(\mathbf{O}(S')\) operates on
the \(S'\)-scheme \(\mathrm{I}_{S'}(\mathcal{M})\), functorially in
\(\mathcal{M}\), \ie the \(S\)-scheme \(\mathrm{I}_S(\mathcal{M})\) is endowed
with the structure of an object with monoid of operators \(\mathbf{O}_S\),
this structure being functorial in \(\mathcal{M}\).

One therefore has a morphism of \(S\)-schemes
\[
\lambda\colon\mathrm{I}_S(\mathcal{M})\times_S\mathbf{O}_S
\longrightarrow\mathrm{I}_S(\mathcal{M}),
\]
satisfying obvious conditions. For every \(S\)-functor \(X\), one obtains by
base change a morphism of \(X\)-functors:
\[
\lambda_X\colon\mathrm{I}_X(\mathcal{M})\times_S\mathbf{O}_S
\longrightarrow\mathrm{I}_X(\mathcal{M})
\]
which makes the \(S\)-functor \(\mathrm{I}_X(\mathcal{M})\) an object with
monoid of operators \(\mathbf{O}(X)\): every element \(a\) of
\(\mathbf{O}(X)=\Hom_S(X,\mathbf{O}_S)\) defines an \(X\)-endomorphism
\[
a^*=\lambda_X\circ\bigl(\mathrm{id}_{\mathrm{I}_X(\mathcal{M})}
\times(a\circ\mathrm{pr}_X)\bigr)
\]
of \(\mathrm{I}_X(\mathcal{M})\); explicitly, if \(x\in X(S')\) and
\(m\in\mathrm{I}_S(\mathcal{M})(S')=\mathrm{I}_{S'}(\mathcal{M})(S')\), then
\(a(x)=a\circ x\) belongs to \(\mathbf{O}(S')\) and one has:
\[
(m,x)\cdot a=(m\cdot a(x),x).
\]
This operation is functorial in \(\mathcal{M}\) and preserves the zero
section \(\varepsilon_{\mathcal{M}}\colon X\to\mathrm{I}_X(\mathcal{M})\),
\ie \(a^*\circ\varepsilon_{\mathcal{M}}=\varepsilon_{\mathcal{M}}\) for every
\(a\in\mathbf{O}(X)\).

Moreover, this operation is ``\emph{functorial in \(X\)}'' in the following
sense: if \(\pi\colon Y\to X\) is a morphism of \(S\)-functors and
\(u\colon\mathbf{O}(X)\to\mathbf{O}(Y)\) the corresponding morphism of rings
(\ie \(u(a)=a\circ\pi\in\mathbf{O}(Y)\) for every
\(a\in\mathbf{O}(X)=\Hom_S(X,\mathbf{O}_S)\)), then the following diagram is
commutative:
\[
\xymatrix{
\mathrm{I}_X(\mathcal{M}) \ar[r]^{a^*} \ar[d]_{\pi}
  & \mathrm{I}_X(\mathcal{M}) \ar[d]^{\pi} \\
\mathrm{I}_Y(\mathcal{M}) \ar[r]^{u(a)^*}
  & \mathrm{I}_Y(\mathcal{M})\,.
}
\]

\numpara{2.2}
Let now \(\mathcal{M}\) and \(\mathcal{N}\) be two quasi-coherent
\(\OO_S\)-modules. The commutative diagram
\[
\xymatrix@C=1.8em@R=1.4em{
& \mathcal{M}\oplus\mathcal{N} \ar[dl] \ar[dr] & \\
\mathcal{M} \ar[dr] & & \mathcal{N} \ar[dl] \\
& 0 &
}
\]
defines a commutative diagram of \(S\)-schemes
\begin{equation}
\vcenter{\xymatrix@C=1.6em@R=1.6em{
& \mathrm{I}_S(\mathcal{M}\oplus\mathcal{N}) & \\
\mathrm{I}_S(\mathcal{M}) \ar[ur]
  & &
\mathrm{I}_S(\mathcal{N}) \ar[ul] \\
& S
  \ar[ul]^{\varepsilon_{\mathcal{M}}}
  \ar[uu]|{\varepsilon_{\mathcal{M}\oplus\mathcal{N}}}
  \ar[ur]_{\varepsilon_{\mathcal{N}}}\,. &
}}
\tag{*}
\end{equation}

\begin{proposition}{2.2}
\label{II.2.2}
For every \(S\)-scheme \(X\), the diagram of functors over \(S\) obtained by
applying the functor \(\Hom_S(-,X)\) to the diagram~\((*)\) is cartesian:
\[
\xymatrix@C=1em@R=1.7em{
& \Hom_S\bigl(\mathrm{I}_S(\mathcal{M}\oplus\mathcal{N}),X\bigr)
  \ar[dl] \ar[dr] & \\
\Hom_S\bigl(\mathrm{I}_S(\mathcal{M}),X\bigr) \ar[dr]
  & &
\Hom_S\bigl(\mathrm{I}_S(\mathcal{N}),X\bigr) \ar[dl] \\
& \Hom_S(S,X)=X\,. &
}
\]
\end{proposition}

One must check that for every \(S'\to S\), the diagram of sets obtained by
taking the value of the functors on \(S'\) is cartesian. As the formation of
\(\mathrm{I}_S(\mathcal{P})\) commutes with extension of the base in the sense
made explicit above, it suffices to do this for \(S'=S\), hence to check that
the following diagram of sets is cartesian:
\[
\xymatrix@C=1.2em@R=1.7em{
& X\bigl(\mathrm{I}_S(\mathcal{M}\oplus\mathcal{N})\bigr)
  \ar[dl] \ar[dd]|{X(\varepsilon_{\mathcal{M}\oplus\mathcal{N}})} \ar[dr] & \\
X\bigl(\mathrm{I}_S(\mathcal{M})\bigr) \ar[dr]_{X(\varepsilon_{\mathcal{M}})}
  & &
X\bigl(\mathrm{I}_S(\mathcal{N})\bigr) \ar[dl]^{X(\varepsilon_{\mathcal{N}})} \\
& X(S)\,. &
}
\]
Now, if \(x\in X(S)\), it follows from SGA~1, III~5.1%
{\renewcommand{\thefootnote}{(11)}\nde{See also the addition~0.1.8 in
Exp.~III.}},
that \(X(\varepsilon_{\mathcal{M}})^{-1}(x)\) is isomorphic, functorially in
\(\mathcal{M}\), to
\[
\Hom_{\OO_S}\bigl(x^*(\Omega^1_{X/S}),\mathcal{M}\bigr),
\]
where \(\Omega^1_{X/S}\) denotes the sheaf of relative differentials of \(X\)
with respect to \(S\). Now this latter functor (in \(\mathcal{M}\)) evidently
transforms a direct sum of \(\OO_S\)-modules into the product of the
corresponding sets, whence the result.
